\paragraph{Notation}
We let $S_i = d_i - a_i$ denote the slack of $J_i$ in the unicast
setting. When requests have non-uniform processing times (or lengths) we
use $\ell_i$ to denote the length of $J_i$. We assume without loss
of generality that $S_i \ge \ell_i$. In the broadcast setting,
$J_{p,i}$ denotes the $i$'th request for page $p$. We assume that the
requests for a page are ordered by arrival time and hence $a_{p,i}
\le a_{p,j}$ for $i < j$. In both settings we use $\Delta$ to denote
the ratio of maximum slack to the minimum slack in a given request
sequence. When pages have non-uniform sizes, $\ell_p$ denotes the size of page $p$.  For an interval $I = [a,b]$ on the real line we use $|I|$
for the length of the interval $(b-a)$. We say a request is alive
 at $t$ if has arrived by $t$ but has not yet finished.
